
%==========================================================================
\section{Discussion : forces et faiblesses}\label{sec:discussion}
%==========================================================================

\subsection{Points forts}

\paragraph{S1 - Quadratique, pas cubique, en $K$.} C'est l'affirmation centrale de
l'article, et la Section~\ref{sec:factors} la confirme dans la forme qui résiste à
l'examen : sur la fenêtre $K=32\rightarrow128$ sur Amazon Movies, une itération eALS
croît de $\GrowthEalsAmz\times$ et une itération ALS de $\GrowthAlsAmz\times$. Comme
la précision de la MF implicite continue de s'améliorer avec $K$, le régime où elle est
la plus performante est précisément le régime où eALS gagne par la plus grande marge
- et le croisement à $K=\FlipKAmz$ se situe en dessous du $K$ que l'on veut utiliser.

\paragraph{S2 - Pas d'inversion de matrice, donc pas de problème de
conditionnement.} Chaque mise à jour de coordonnée est une division scalaire. ALS
doit inverser $Q^{\top}W^uQ+\lambda I$ une fois par utilisateur ; à $K=256$ avec un
$Q$ mal conditionné, c'est là que vivent les ennuis numériques, et c'est aussi
pourquoi MLlib a besoin d'un chemin Cholesky avec un plancher de régularisation. eALS
n'a pas un tel objet.

Une réserve que l'article ne soulève pas, et que nous pouvons préciser. Le
dénominateur de~\eqref{eq:pu} est
$\sum_{i\in\mathcal{R}_u}(w_{ui}-c_i)q_{if}^2+s^q_{ff}+\lambda$, dans lequel
$s^q_{ff}=\sum_i c_iq_{if}^2>0$ et $\lambda>0$, mais la première somme n'est
\emph{pas} de signe défini : rien dans le modèle n'interdit $c_i>w_{ui}$. Comme
$c_i \approx c_0/N$ en moyenne, cela se produit sur les petits catalogues - sur
ml-100k avec le $c_0=512$ de l'article et $N=1\,152$, l'item le plus populaire atteint
$c_i=1,07>1=w_{ui}$. Nous n'avons jamais observé d'objectif non monotone
(Figure~\ref{fig:conv}), car $s^q_{ff}$ domine en pratique, mais la garantie est
empirique plutôt que structurelle, et c'est une raison de faire varier $c_0$ avec $N$
plutôt que de transposer une valeur d'un jeu de données à l'autre.

\paragraph{S3 - Un parallélisme exact.} La
Section~\ref{sec:correctness} mesure un écart maximal de $\CorrWorst$ entre la
référence NumPy séquentielle et les exécutions Spark à 1, 2, 4 et 8 blocs, et
l'exécution à un bloc est identique au bit près. Rien n'a besoin d'être ajusté,
aucune obsolescence (\emph{staleness}) n'a besoin d'être tolérée, aucun taux
d'apprentissage n'a besoin d'être re-recherché quand le nombre de workers change -
ce que la SGD distribuée ne peut précisément pas promettre.
Opérationnellement cela vaut plus qu'il n'y paraît : cela signifie que le nombre de
partitions est un pur bouton de performance, et qu'une expérience de scalabilité ne
peut pas changer silencieusement le modèle qu'elle mesure.

\paragraph{S4 - Pas de taux d'apprentissage.} Chaque mise à jour est le minimiseur
exact de l'objectif le long d'une coordonnée. Comparé à RCD~\cite{devooght2015rcd},
qui a besoin d'une recherche linéaire à chaque pas, ou à BPR~\cite{rendle2009bpr}, dont
la précision est visiblement fonction de la taille du pas, cela retire une dimension
entière du budget d'ajustement. Notre propre expérience le confirme : aucune
configuration que nous ayons exécutée n'a jamais divergé, et l'objectif décroît de
façon monotone à chaque itération sur les deux jeux de données
(Figure~\ref{fig:conv}).

\subsection{Points faibles}

\paragraph{W1 - Le balayage de coordonnées est intrinsèquement séquentiel à
l'intérieur d'une ligne.} La boucle en $K$ est du Gauss--Seidel : $p_{u,f+1}$ doit
voir le $p_{uf}$ mis à jour. Le parallélisme est donc limité à $\min(M,B)$ dans
l'étape utilisateur et $\min(N,B)$ dans l'étape item, et le travail par bloc contient
toujours une boucle Python de longueur $K$. C'est cette boucle qui a obligé
l'implémentation à être vectorisée \emph{entre} utilisateurs : c'est le seul axe
restant. Dans un contexte à peu d'utilisateurs et à $K$ énorme - l'opposé d'un
système de recommandation - eALS se paralléliserait mal.

\paragraph{W2 - Diffuser $Q$ ne passe pas à l'échelle des catalogues réels.} La
conception RDD exige $P$ et $Q$ sur le driver et une copie désérialisée de chacun par
worker Python. À $N=10^6$ items et $K=128$ cela fait 1\,Go par worker. C'est la limite
structurelle de la conception retenue en Section~\ref{sec:solution}, et elle est
assumée : l'alternative - faire transiter les facteurs par des jointures avec
shuffle - déplace $O(|\mathcal{R}|K)$ au lieu de $O((M+N)K)$, soit un facteur
$\approx 300$ sur ml-1m, et se paie à chaque itération. Il n'y a pas de repas gratuit
ici : à une certaine taille de catalogue, la conception fondée sur les jointures gagne
simplement parce que la conception broadcast cesse de tenir en mémoire, et la bonne
réponse industrielle n'est probablement ni l'une ni l'autre, mais le schéma à
partitionnement par blocs de MLlib, où chaque ligne de facteur n'est envoyée qu'aux
blocs qui en ont besoin.

\paragraph{W3 - Le schéma de pondération exige une recherche bidimensionnelle
coûteuse.} $c_0$ et $\alpha$ interagissent, les propres optima de l'article diffèrent
entre ses deux jeux de données ($c_0=512,\alpha=0,4$ sur Yelp ; $c_0=64,\alpha=0,5$
sur Amazon), et chaque point de la grille est un entraînement complet - une
recherche bidimensionnelle sérieuse représente donc des dizaines d'entraînements. Nous
avons repris les valeurs de l'article plutôt que de les rechercher, ce qui est
précisément la raison pour laquelle nos chiffres de précision ne doivent pas être lus
comme un optimum. Pire, la sensibilité n'est pas symétrique : un $c_0$ trop petit
détruit la précision, un $c_0$ trop grand la dégrade doucement, donc une grille
grossière peut facilement tomber sur un plateau et déclarer victoire.

\paragraph{W4 - Un protocole d'évaluation indulgent, en particulier envers la
pondération par popularité.} Trois problèmes distincts :
\begin{itemize}\itemsep2pt
\item \emph{Le leave-one-out par horodatage n'est pas un découpage réaliste.} Il
laisse fuir le futur : l'ensemble d'entraînement de l'utilisateur $u$ contient des
interactions survenues après l'interaction de test de l'utilisateur $v$. L'article le
dit lui-même et ajoute pour cette raison un protocole en ligne fondé sur une coupure
chronologique globale, que nous n'avons pas reproduit ; tous nos chiffres de précision
héritent donc de ce biais.
\item \emph{Des négatifs échantillonnés rendraient les chiffres incomparables.} Nous
classons contre le catalogue entier précisément parce qu'échantillonner 100 négatifs~\cite{krichene2020sampled}
est connu pour changer non seulement les valeurs mais l'\emph{ordre} relatif des
systèmes comparés (Krichene et Rendle, KDD 2020). Cela coûte $O(MNK)$ par évaluation et c'est
la principale raison pour laquelle notre évaluation est distribuée.
\item \emph{HR et NDCG ne mesurent rien sur la diversité ou la couverture du
catalogue.} C'est la critique que nous jugeons la plus importante, car elle pointe
vers la contribution même de l'article. Pondérer les données manquantes par la
popularité indique au modèle qu'un manque sur un blockbuster est une forte preuve de
désintérêt - mais le même mécanisme pousse les items impopulaires vers un score
prédit bas pour tout le monde, car il est bon marché de se tromper sur eux. Une
métrique qui demande seulement « l'item retenu était-il dans le top 100 » ne peut pas
voir la concentration qui en résulte, et les items retenus sont eux-mêmes distribués
selon la popularité, donc la métrique est biaisée positivement en faveur de tout ce
qui amplifie la popularité. Nous mesurons un vrai
gain à partir de $\alpha>0$ (Section~\ref{sec:quality}), et nous croyons qu'il est réel ;
nous croyons aussi qu'une évaluation rapportant la couverture dans l'espace des items
montrerait qu'une partie de ce gain est un biais de popularité plutôt qu'une
meilleure modélisation des préférences. Mesurer cela est l'extension la plus utile de
ce travail.
\end{itemize}

\paragraph{W5 - eALS optimise une perte de régression, pas une perte de
classement.} L'objectif est l'erreur quadratique sur une matrice $0/1$. C'est un bon
objectif de \emph{rappel}, ce que montrent exactement nos résultats - eALS est le
plus fort en HR@100 - et un objectif de \emph{précision} médiocre. BPR, qui
optimise directement l'ordre par paires, est le concurrent naturel aux petites
coupures ; cette asymétrie entre rappel et précision est un résultat classique de la
littérature sur la recommandation top-$N$~\cite{cremonesi2010topn}. Nous n'avons pas implémenté BPR (voir plus
bas) et ne pouvons donc pas le quantifier sur nos données ; l'article rapporte le même
schéma.

Un point connexe et plus inconfortable : Rendle et al. ont
revisité iALS - la référence même à laquelle eALS est comparé - sur des
benchmarks standards et ont trouvé qu'une fois sa régularisation et son poids sur les
données non observées correctement ajustés, il rivalise avec des méthodes bien plus
récentes. Notre propre référence ALS est celle de MLlib, avec le $\lambda=0,01$ de
l'article et l'$\alpha$ par défaut de MLlib ; nous ne l'avons \emph{pas} ajustée. Une
partie de l'écart de $\EalsOverAls$\,\% que nous mesurons est donc attribuable à un
effort d'ajustement inégal, et une lecture honnête du Tableau~\ref{tab:quality}
devrait le traiter comme une borne supérieure de l'avantage d'eALS sur un iALS bien
ajusté.

\subsection{Limites de validité des mesures}

Nous les listons car la plupart sont structurelles et ne peuvent pas être corrigées
sur un portable.

\begin{enumerate}\itemsep2pt
\item \textbf{\texttt{local[$p$]} n'est pas un cluster.} Il n'y a pas de réseau, pas
de sérialisation entre machines, pas de tâche à la traîne due à un nœud lent, pas de
tolérance aux pannes exercée. Nos chiffres \emph{surestiment} donc la performance de
la conception RDD sur un vrai cluster, où diffuser $Q$ coûte de la bande passante
réseau et non une copie mémoire, et \emph{sous-estiment} d'autant la valeur de
l'alternative co-partitionnée écartée en Section~\ref{sec:solution}.
\item \textbf{La bande passante mémoire partagée est la vraie limite, pas le nombre
de cœurs.} Le noyau est une séquence de gathers et de sommes par segment sur des
tableaux bien plus grands que le L2. Passer de 1 à 8 cœurs multiplie la demande sur un
seul contrôleur mémoire, ce qui explique une grande partie de l'écart entre
l'accélération mesurée $\SpeedupBig\times$ et l'idéal $8\times$. Un cluster de 8
machines à un cœur paraîtrait meilleur qu'une seule machine à 8 cœurs ici.
\item \textbf{La dérive thermique est mesurable, et nous l'avons attrapée.} Deux
exécutions du même balayage en $K$ sur ml-1m, à une heure d'intervalle dans une
campagne continue, différaient jusqu'à un facteur deux sur les configurations les
plus lourdes : la seconde exécution suivait une mesure ALS bien plus lourde et la
machine était chaude. Un portable sous charge soutenue à 8 cœurs ne maintient pas une
horloge constante. Nos parades sont méthodologiques plutôt que physiques -
répétitions entrelacées entre configurations plutôt que regroupées, et un temps de
repos entre configurations (\texttt{--interleave}, \texttt{--cooldown}) - et les
figures montrent des barres min--max sur 9 échantillons par itération afin que
l'étalement résiduel soit visible. Sur une machine dédiée à horloge fixe, ces
précautions seraient inutiles ; sur celle-ci, elles font la différence entre une
mesure et une anecdote.
\item \textbf{Deux familles de jeux de données, un seul domaine chacune.} Amazon
Movies et MovieLens sont tous deux de nature cinématographique, tous deux
modérément denses après filtrage 10-core. Rien ici ne dit comment eALS se comporte
sur un catalogue de $10^7$ items avec une moyenne de deux interactions chacun.
\item \textbf{Le notebook est le seul test de bout en bout.} Il n'y a pas de suite de
tests unitaires ; l'exactitude est garantie par \texttt{check\_correctness.py} (qui
compare bien à une référence par force brute) et par le fait que le notebook
s'exécute jusqu'au bout. Cela a suffi à attraper de vrais bugs - notamment un
\texttt{PYTHONPATH} manquant sur les workers - mais ce n'est pas un substitut à des
tests.
\item \textbf{Tout ce qui figurait dans l'énoncé n'a pas été mesuré.} Nous n'avons
pas implémenté BPR ni RCD : les deux sont
des références pertinentes dans l'article, mais chacune est un projet en soi et le
budget est allé à l'étude de scalabilité que l'énoncé pondère le plus. Nous le disons
plutôt que de rapporter des chiffres que nous n'avons pas produits. Nous n'avons pas
non plus utilisé Yelp, pour la raison donnée en Section~\ref{sec:dataset}.
\end{enumerate}
